\subsection{层的例子}

设 B 为一个拓扑空间。

\subitemhead{映射的层}

设 X 为一个集合。对 B 的每个开集 U，用 \(\mathcal{F}(\mathrm{U};\mathrm{X})\) 表示从 U 到 X 中的映射所成的集 (E, II, p. 31)。对 B 的满足 \(\mathrm{U}\subset \mathrm{V}\) 的每对开集 (U, V)，令 \(r_{\mathrm{UV}}\colon \mathcal{F}(\mathrm{V};\mathrm{X})\to\mathcal{F}(\mathrm{U};\mathrm{X})\) 为限制映射 \(f\mapsto f|\mathrm{U}\) 。显然，对 \((\mathcal{F}(\mathrm{U};\mathrm{X}),r_{\mathrm{UV}})\) 是 B 上的一个层。它称为 B 上的以 X 为值的映射的层，记作 \(\underline{\mathcal{F}} (\mathrm{B};\mathrm{X})\) 。

\subitemhead{连续映射的层}

设 X 为一个拓扑空间。对 B 的每个开集 U，令 \(\mathcal{C}(\mathrm{U};\mathrm{X})\) 为从 U 到 X 中的连续映射所成的集。那么，\((\mathcal{C}(\mathrm{U};\mathrm{X}))\) 是层 \(\underline{\mathcal{F}} (\mathrm{B};\mathrm{X})\) 的一个子层（由 TG, I, p. 19 的 prop. 4）。这样得到的层记作 \(\underline{\mathcal{C}} (\mathrm{B};\mathrm{X})\) ，称为 B 上的以 X 为值的连续映射的层。在 X 赋有离散拓扑的特殊情形，层 \(\underline{\mathcal{C}} (\mathrm{B};\mathrm{X})\) 取名为 B 上的以 X 为值的局部常值函数的层。

\subitemhead{连续截面的层}

设 E 为一个拓扑空间，并设 \(p \colon \mathrm{E} \to \mathrm{B}\) 为一个连续映射。对 B 的每个开集 U，我们用 \(\mathcal{S}(\mathrm{U};p)\) （当不致混淆时也写作 \(\mathcal{S}(\mathrm{U};\mathrm{E})\) ）表示 \(p\) 在 U 上的连续截面所成的集。族 \((\mathcal{S}(\mathrm{U};p))\) 是层 \(\underline{\mathcal{C}}(\mathrm{B};\mathrm{E})\) 的一个子层。这样得到的层记作 \(\underline{\mathcal{L}}(\mathrm{B};\mathrm{E})\) 或简称 \(\underline{\mathcal{S}}(\mathrm{E})\) ，称为 B-空间 \((\mathrm{E},p)\) 的连续截面在 B 上的层。我们将在下文第 6 号中看到，B 上的每个层都同构于某个平展 B-空间的连续截面在 B 上的层。

\subitemhead{B-态射的层}

设 \((\mathrm{E},p)\) 与 \((\mathrm{E}^{\prime},p^{\prime})\) 为 B-空间。对 E 的每个开集 U，我们用 \(\mathcal{C}_{\mathrm{B}}(\mathrm{U};\mathrm{E}^{\prime})\) 表示从 \((\mathrm{U},p|\mathrm{U})\) 到 \((\mathrm{E}^{\prime},p^{\prime})\) 中的 B-态射所成的集。族 \((\mathcal{C}_{\mathrm{B}}(\mathrm{U};\mathrm{E}^{\prime}))\) 是层 \(\underline{\mathcal{C}}(\mathrm{E};\mathrm{E}^{\prime})\) 的一个子层。这样得到的层记作 \(\underline{\mathcal{C}}_{\mathrm{B}}(\mathrm{E};\mathrm{E}^{\prime})\) ，称为以 \((\mathrm{E}^{\prime},p^{\prime})\) 为值的 B-态射在 E 上的层。当 \((\mathrm{E},p)\) 等于 \((\mathrm{B},\mathrm{Id}_{\mathrm{B}})\) 时，这个层就是例 3 中的层 \(\underline{\mathcal{L}}(\mathrm{B};\mathrm{E})\) 。

对 B 的每个开集 U，令 \(\mathcal{M}(\mathrm{U})\) 为从 \(p^{-1}(\mathrm{U})\) 到 \((p')^{-1}(\mathrm{U})\) 中的 U-态射所成的集。对 B 的每对开子集 \((\mathrm{U},\mathrm{V})\) （满足 \(\mathrm{U}\subset \mathrm{V}\) 者），令 \(m_{\mathrm{UV}}\colon \mathcal{M}(\mathrm{V})\to \mathcal{M}(\mathrm{U})\) 为这样的映射：它把每个 V-态射 \(f\colon p^{-1}(\mathrm{V})\to (p')^{-1}(\mathrm{V})\) 映为从 \(p^{-1}(\mathrm{U})\) 到 \((p')^{-1}(\mathrm{U})\) 中的 U-态射

（由 \(f\) 过渡到子集而得者）。那么 \((\mathcal{M}(\mathrm{U}), m_{\mathrm{UV}})\) 是 B 上的一个层。我们把它记作 \(\underline{\mathcal{M}or}_{\mathrm{B}}(\mathrm{E}; \mathrm{E}')\) ，并称它为从 \((\mathrm{E}, p)\) 到 \((\mathrm{E}', p')\) 中的 B-态射在 B 上的层。

对 B 的每个开集 U，令 \(\mathcal{I}s(\mathrm{U})\) 为 \(\mathcal{M}(\mathrm{U})\) 中由从 \(p^{-1}(\mathrm{U})\) 到 \((p')^{-1}(\mathrm{U})\) 的 U-同构所组成的子集。族 \((\mathcal{I}s(\mathrm{U}))\) 是态射层 \(\underline{\mathcal{Mor}}_{\mathrm{B}}(\mathrm{E};\mathrm{E}^{\prime})\) （从 \((\mathrm{E},p)\) 到 \((\mathrm{E}^{\prime},p^{\prime})\) 中的态射的）的一个子层。这样得到的层记作 \(\underline{\mathcal{Isom}}_{\mathrm{B}}(\mathrm{E};\mathrm{E}^{\prime})\) ，称为从 \((\mathrm{E},p)\) 到 \((\mathrm{E}^{\prime},p^{\prime})\) 上的 B-同构在 B 上的层。

\subitemhead{\(\mathbf{C}^r\) 类映射的层}

设 X 与 Y 为域 K 上的 \(C^{r}\) 类簇（关于 K 与 r 的约定取 VAR, R 中的约定）。对 X 的每个开子集 U，令 \(\mathcal{C}^{r}(\mathrm{U};\mathrm{Y})\) 为从 U 到 Y 中的 \(C^{r}\) 类映射所成的集。族 \((\mathcal{C}^{r}(\mathrm{U};\mathrm{Y}))\) 是层 \(\underline{\mathcal{C}}(\mathrm{X};\mathrm{Y})\) 的一个子层。这样得到的层记作 \(\underline{\mathcal{C}}^{r}(\mathrm{X};\mathrm{Y})\) ，称为 X 上的以 Y 为值的 \(C^{r}\) 类映射的层（参看 VAR, R, 5.4.2）。

\subitemhead{子空间的层}

若 U 与 V 是 B 的满足 \(U \subset V\) 的开子集，用 \(i_{\mathrm{UV}}: \mathfrak{P}(V) \to \mathfrak{P}(U)\) 表示把 V 的每个子集 A 映为子集 \(A \cap U\) 的映射。对 \((\mathfrak{P}(U), i_{\mathrm{UV}})\) 是一个层，称为 B 的子空间的层，记作 \(\underline{\mathfrak{P}}(B)\) 。事实上，若用 X 表示集 \(\{0;1\}\)，则把 U 的每个子集 A 映为其特征函数 \(\varphi_{A}^{U}: U \to X\) 的映射是从 \(\mathfrak{P}(U)\) 到 \(\mathcal{F}(U;X)\) 上的一个双射 (E, III, p. 38)；此外，若 U 与 V 是满足 \(U \subset V\) 的开集，则对 V 的每个子集 A，\(\varphi_{A \cap U}^{U}\) 是 \(\varphi_{A}^{V}\) 在 U 上的限制，从而 \(\underline{\mathfrak{P}}(B)\) 可等同于 B 上的以 X 为值的映射的层。

对 B 的每个开集 U，设 \(\mathcal{L}(\mathrm{U})\) 为 \(\mathfrak{P}(\mathrm{U})\) 的一个子集。要使 \((\mathcal{L}(\mathrm{U}))\) 是 \(\underline{\mathfrak{P}}(\mathrm{B})\) 的子层，必须且只需下列条件被满足：

(F') 设 \((\mathrm{U}_{i})_{i\in\mathrm{I}}\) 为 B 的开集的一个族，U 为其并，并设 A 为

U 的一个子集；要使 A 属于 \(\mathcal{L}(\mathrm{U})\) ，必须且

只需，对 I 中的每个 i，\(A \cap \mathrm{U}_{i}\) 属于 \(\mathcal{L}(\mathrm{U}_{i})\)。

例如，若 \(\mathcal{L}(\mathrm{U})\) 是 U 的闭子集所成的集，则条件 \((\mathrm{F}')\) 满足。

\subitemhead{层的积}

设 \(\mathcal{B}\) 为 B 的拓扑的一个基，并设 I 为一个集合。对每个 \(i \in I\) ，设 \(\mathcal{F}_{i} = (\mathcal{F}_{i}(U), f_{i,\mathrm{UV}})\) 为 B 上的、相对于基 \(\mathcal{B}\) 的一个预层。对每个开集 \(U \in \mathcal{B}\) ，令 \(\mathcal{F}(U) = \prod_{i \in I} \mathcal{F}_{i}(U)\) ，并对 \(\mathcal{B}\) 中满足 U ⊂ V 的每对元素 (U, V)，用 \(f_{\mathrm{UV}}\) 表示映射 ( \(f_{i,\mathrm{UV}}\) ) \(_{i\in\mathrm{I}}\) ：\(\mathcal{F}(\mathrm{V}) \to \mathcal{F}(\mathrm{U})\) 。那么 ( \(\mathcal{F}(\mathrm{U}), f_{\mathrm{UV}} \) ) 是 B 上的、相对于 \(\mathcal{B}\) 的一个预层；它称为族 ( \(\mathcal{F}_i\) ) 的积预层，记作 \(\prod_{i\in\mathrm{I}}\mathcal{F}_i\) 。若对每个 \(i \in \mathrm{I}\) ，\(\mathcal{F}_i\) 都是层，则它是一个层。

